\begin{abstract}
Private Set Union (PSU) enables two parties to compute the union of their
private sets. While recent IBLT-based protocols achieve PSU performance
competitive with Private Set Intersection, they are \emph{static}: computing the
union on updated sets requires re-running the full protocol from scratch.

We present \textbf{FUPSU}, the first forward-private updatable PSU protocol.
FUPSU exploits the linearity of IBLT encoding to support incremental
(delta-based) set updates, achieving per-epoch communication and computation
$O(|\Delta|)$. We identify and resolve a fundamental barrier to IBLT-based
incremental updates: \emph{deletion cascade}, where removing even a few elements
triggers global propagation through the IBLT structure, erasing the benefit of
incremental processing. Our key insight---\textbf{OT-masked deletion}---offloads
deletion membership checks to Oblivious Transfer, preventing shared-element
deletions from entering the IBLT cascade path entirely. This restores true
$O(|\Delta|)$ scaling, yielding $29$--$23{,}756\times$ per-epoch cell scan
reduction over static full re-encoding, across update sizes $|\Delta|=64$--$4096$.
Forward privacy is obtained via a PRF-based key evolution chain with immediate
erasure, adding negligible overhead. We also prove a \emph{leakage lower bound}:
any updatable PSU protocol with sublinear communication must leak individual
add and delete counts---our protocol matches this bound and is therefore
leakage-optimal.

We further show that a single execution of the UnionPeel sub-protocol
simultaneously recovers the union, the intersection, and their
cardinalities---yielding a \emph{unified} updatable private set operation
framework (updPSO) with zero additional communication or computation.
We implement FUPSU end-to-end in C++ and benchmark across set sizes $2^{14}$
to $2^{20}$ with $100\%$ recovery across all configurations. All code is
open source.
\end{abstract}

\section{Introduction}
% ============================================================

Private Set Union (PSU)~\cite{KRT19,JLZ20,ZCL+23,Tu+25} allows two parties to
compute the union $X \cup Y$ of their private sets while revealing nothing
beyond the union itself. With applications in privacy-preserving identity
management, collaborative cyber-risk assessment, private database joins, and
federated learning~\cite{DCW13,IKN+20,KRS+19}, PSU is a fundamental secure
computation primitive alongside Private Set Intersection (PSI).

Historically, PSU has been significantly slower than PSI---often by $30\times$
or more~\cite{PT26,KRT19}. The recent breakthrough of Piske and
Trieu~\cite{PT26} (Eurocrypt 2026) closed this gap by introducing the first
IBLT-based PSU protocol. Their key insight, \emph{union peelability}, enables
peeling the union of two IBLTs without materializing the combined table, using
oblivious transfer (OT) and an oblivious PRF (OPRF). The resulting protocol
computes PSU on $2^{20}$ elements in under 3 seconds, matching state-of-the-art
PSI performance.

\textbf{The static limitation.} Piske and Trieu's protocol, like all prior PSU
constructions, is \emph{static}: both parties must encode their entire sets into
fresh IBLTs for each computation. In many real-world applications, however,
sets evolve incrementally---new elements are added, old ones expire---in applications
including threat intelligence sharing, federated learning with evolving client
datasets, and private database join maintenance under incremental updates.
Re-running the full protocol on each update wastes resources proportional
to the total set size, even when only a handful of elements change.

\textbf{The forward privacy requirement.} The multi-epoch nature of updatable
protocols introduces a new security concern: \emph{forward privacy}. If a party
is compromised at epoch $\tau$, the adversary should not be able to recover the
party's historical set elements from epochs $t < \tau$, even having observed all
protocol messages. This is the privacy analogue of forward secrecy in
Signal~\cite{Signal}, and is critical when sets contain sensitive data with
limited lifetimes. Privacy regulations (GDPR, CCPA) mandate data retention
limits; if historical elements remain decryptable after key compromise, the
system violates the principle that deleted data should be irrecoverable.

\textbf{Existing updatable work.} Several updatable PSI protocols
exist~\cite{BMX22,BMS+24,LMR+26,ABF+25,LTS+26}, with FUPSI~\cite{WZC+24}
achieving forward privacy via keyword PIR. \textbf{No prior work} addresses updatable PSU,
let alone forward-private updatable PSU.

\textbf{Our contributions.} We initiate the study of updatable PSU with forward
privacy:

\begin{enumerate}
\item \textbf{The FUPSU protocol} (Section~\ref{sec:protocol}). The first
forward-private updatable PSU protocol. Built on three mechanisms: (i)~delta IBLT
encoding with $O(|\Delta|)$ per-epoch cell operations;
size $n$; (ii)~\textbf{OT-masked deletion}, which offloads deletion membership
checks to Oblivious Transfer, preventing shared-element deletions from entering
the IBLT cascade path; (iii)~split-path incremental peeling with cascade bounded
by $2k|\Delta|$; (iv)~PRF-based key evolution for forward privacy. UnionPeel
simultaneously recovers the intersection and cardinalities, yielding a unified
updPSO framework.

\item \textbf{Leakage lower bound} (Section~\ref{sec:protocol}, Theorem~\ref{thm:lower}). We prove that
\emph{any} $1$-updatable PSU protocol must reveal the individual add and delete
counts for each party---an inherent efficiency-leakage trade-off. Any protocol
hiding these counts requires $\Omega(N\sigma)$ communication, defeating
updatability. FUPSU's OT queries and delta IBLT sizes reveal these four counts,
matching the lower bound and making it \textbf{leakage-optimal}.

\item \textbf{Implementation and evaluation} (Section~\ref{sec:eval}).
We implement FUPSU end-to-end in C++ (IBLT with $k{=}5$, OT-masked deletion,
MiniPeel for additions). Benchmarks on Intel Xeon E5-2630 v4 across set sizes
$2^{14}$--$2^{20}$ and update sizes $|\Delta|=64$--$4096$ achieve
$29$--$23{,}756\times$ per-epoch cell scan reduction over static full
re-encoding, with $100\%$ recovery across all configurations. We identify
the deletion cascade barrier and quantify how OT-masking eliminates it.
\end{enumerate}

\smallskip\noindent
On the security front, we prove FUPSU UC-realizes~\cite{Can01}
$\cF_{\mathsf{updPSU}}$ in the $(\cF_{\mathsf{OPRF}},
\cF_{\mathsf{OT}})$-hybrid model against semi-honest adaptive adversaries,
with a modular proof that separates per-epoch PSU security from cross-epoch
key management, enabling independent upgrades. We introduce FP-IND, a
standalone game-based definition that cleanly captures the forward privacy
guarantee of the PRF chain. Because PSU outputs the union, FUPSU is
structurally unaffected by adaptive input-selection attacks
on UPSI~\cite{LMR+26}, requiring no protocol modification.

\smallskip
\noindent\textbf{Comparison with prior work.}
Table~\ref{tab:intro-compare} situates FUPSU among prior PSU and
updatable protocols.

\begin{table*}[t]
\centering
\footnotesize
\begin{tabular}{|l|l|c|c|c|c|c|}
\hline
\textbf{Protocol} & \textbf{Type} & \textbf{Updatable} & \textbf{Fwd.\ Priv.} & \textbf{UC} & \textbf{Adaptive} & \textbf{Complexity} \\
\hline\noalign{\vspace{-1pt}}\hline
KRT19~\cite{KRT19}   & PSU & No  & No  & Yes & No  & $O(n)$ \\
ZCL+23~\cite{ZCL+23} & PSU & No  & No  & Yes & No  & $O(n)$ \\
Jia+24~\cite{Jia+24} & PSU & No  & No  & Enhanced & No  & $O(n \log n)$ \\
Tu+25~\cite{Tu+25}   & PSU & No  & No  & Enhanced & No  & $O(n)$ \\
PT26~\cite{PT26}     & PSU & No  & No  & Yes & No  & $O(n)$ \\
FUPSI~\cite{WZC+24}  & PSI & Yes & Yes & No  & No  & $O(\log n)$ \\
LMR+26~\cite{LMR+26} & PSI & Yes & No  & Yes & Yes & $O(|\Delta|)$ \\
\hline\noalign{\vspace{-1pt}}\hline
\textbf{Ours} & PSU & Yes & Yes & Yes & Yes & $O(|\Delta|)$ \\
\hline
\end{tabular}

\smallskip
\caption{Comparison with prior PSU and updatable protocols.
Enhanced = UC-secure under the enhanced PSU functionality (no within-execution leakage).}
\label{tab:intro-compare}
\end{table*}

% ============================================================

\subsection{Technical Overview}

We now give an informal walkthrough of our key technical ideas. The formal protocol specification appears in Section~\ref{sec:protocol}, with full proofs in Sections~\ref{sec:security} and the Appendix.

\subsubsection{The Deletion Cascade Barrier}\label{sec:cascade-barrier}

A natural incremental approach applies delta IBLT encoding: deletions are
subtracted from IBLT cells, and UnionPeel starts from the $k|\Delta^-|$
modified cells. This directly realizes the Cascade Bound Lemma
(Lemma~\ref{lem:cascade}), which proves $\leq 2k|\Delta^-|$ cells are
examined. However, \emph{in practice}, the cascade does not stop there.

The issue is structural. After subtracting deleted elements, new singletons
emerge in the IBLT. Peeling these singletons exposes further singletons---a
chain reaction that propagates through \emph{all remaining shared elements}
in the IBLT, not just those related to deletions. Empirically, a single
deletion sweep peels $\approx 70$--$80\%$ of all shared elements, examining
$\Theta(|S|)$ cells instead of $O(|\Delta|)$. This reduces the incremental
benefit to only $6$--$10\times$ over static re-encoding, far below the
theoretical $n/|\Delta|$ potential.

\subsubsection{OT-Masked Deletion}

Our key insight is that \emph{most deletions do not require IBLT
modification at all}. When a party deletes an element $x$, the protocol
only needs to know whether $x$ is (i)~unique to the deleting party,
(ii)~shared but retained by the other party, or (iii)~deleted by both.
Using a 2-bit OT-based membership check---already supported by the
$\cF_{\mathsf{OPRF}}$ and $\cF_{\mathsf{OT}}$ functionalities used by
UnionPeel---the parties classify each deletion into one of three cases:

\begin{itemize}
\item \textbf{Case~1 (unique):} The other party does not hold $x$. The element
  is removed from the union via local update only---no IBLT modification.
\item \textbf{Case~2 (one-sided shared):} The other party holds $x$ and does
  not delete it. $x$ remains in the union. The deleting party subtracts $x$
  from its IBLT ($k$ cell updates, no cascade).
\item \textbf{Case~3 (two-sided shared):} Both parties independently delete
  $x$. This is the only case requiring IBLT subtraction ($\ominus$~on both
  parties' IBLTs).
\end{itemize}

In practice, Case~3 is rare ,
because the probability of both parties independently selecting the same
shared element for deletion is low. The remaining deletions are
handled purely locally (Case~1) or via OT without IBLT touch (Case~2). This
\emph{eliminates deletion cascade entirely}---per-epoch cell scanning is
dominated by the Mini UnionPeel for additions, at $O(|\Delta^+|)$ cells.
The resulting per-epoch cost is true $O(|\Delta|)$, restoring the full
$n/|\Delta|$ reduction factor.